\chapter{The syntax in brief}\label{app-syntax}

The Revised Report defines the syntax of Algol 68 with a two-level
grammar, which also expresses the type rules. This appendix gives a
simplified, context-free summary in the extended Backus--Naur form that
most language references use today. It describes the shape of programs;
whether a program is correct also depends on the types, which the grammar
below does not check. The names of the nonterminals are the terms of
this book, with the Report's term in a comment where it differs.

Notation: \texttt{\{ x \}} means zero or more repetitions of \texttt{x},
\texttt{[ x ]} an optional \texttt{x}, \texttt{x | y} a choice; terminal
symbols are in quotes.

\begin{lstlisting}[language={},basicstyle=\ttfamily\footnotesize]
program        = [ access ] enclosed .
access         = "access" module_name { "," module_name } .

block          = "begin" sequence "end" | "(" sequence ")" .
                                          (* closed clause *)
sequence       = { declaration ";" } part { exit_part } .
                                          (* serial clause *)
part           = { label ":" } expression
                 { ";" ( declaration | { label ":" } expression ) } .
exit_part      = "exit" label ":" part .

declaration    = mode_decl | prio_decl | op_decl | constant_decl
               | variable_decl | proc_decl .
mode_decl      = "mode" Type_name "=" type { "," Type_name "=" type } .
prio_decl      = "prio" operator "=" digit { "," operator "=" digit } .
op_decl        = "op" operator "=" function_literal
                 { "," operator "=" function_literal } .
constant_decl  = formal_type name "=" expression
                 { "," name "=" expression } .  (* identity declaration *)
variable_decl  = [ "loc" | "heap" ] actual_type name [ ":=" expression ]
                 { "," name [ ":=" expression ] } .
proc_decl      = "proc" name "=" function_literal
                 { "," name "=" function_literal } .

type           = "int" | "real" | "bool" | "char" | "string" | "compl"
               | "bits" | "bytes" | "sema" | "format" | "void"
               | ( "long" | "short" ) { "long" | "short" } type
               | "ref" type
               | "[" [ bounds { "," bounds } ] "]" type
               | "flex" "[" [ bounds { "," bounds } ] "]" type
               | "struct" "(" type names { "," type names } ")"
               | "union" "(" type { "," type } ")"
               | "proc" [ "(" type { "," type } ")" ] type
               | Type_name .
bounds         = [ expression ":" ] expression .
                    (* formal types have no bounds, actual ones have *)

expression     = assignment | identity | function_literal | jump
               | "skip" | "nil" | formula
               | tertiary ( "andth" | "orel" ) tertiary .  (* unit *)
assignment     = tertiary ":=" expression .
identity       = tertiary ( "is" | "isnt" | ":=:" | ":/=:" ) tertiary .
function_literal = [ "(" type names { "," type names } ")" ] type ":"
                   expression .                 (* routine text *)
jump           = [ "goto" | "go" "to" ] label .
formula        = [ formula operator ] operand .
                    (* binary operators by precedence 1 to 9 *)
operand        = { unary_operator } secondary .
secondary      = primary | selection | generator .
selection      = field "of" secondary .
generator      = ( "loc" | "heap" ) actual_type .
primary        = name | literal | cast | call | slice | enclosed .
cast           = type enclosed .
call           = primary "(" expression { "," expression } ")" .
slice          = primary "[" index { "," index } "]" .
index          = expression | [ expression ] ":" [ expression ]
                 [ "@" expression ] | "@" expression .

enclosed       = block | collateral | parallel | if_expr | case_expr
               | conformity | loop .
collateral     = "(" expression "," expression { "," expression } ")"
               | "(" ")" .              (* also array and record literals *)
parallel       = "par" "(" expression "," expression { "," expression } ")" .
if_expr        = "if" sequence "then" sequence
                 { "elif" sequence "then" sequence }
                 [ "else" sequence ] "fi" .
case_expr      = "case" sequence "in" sequence { "," sequence }
                 [ "ouse" ... | "out" sequence ] "esac" .
conformity     = "case" sequence "in"
                 "(" type [ name ] ")" ":" sequence
                 { "," "(" type [ name ] ")" ":" sequence }
                 [ "out" sequence ] "esac" .
loop           = [ "for" name ] [ "from" expression ] [ "by" expression ]
                 [ "to" expression ] [ "while" sequence ]
                 "do" sequence "od" .
\end{lstlisting}

The brief forms \code{( c | a | b )} for \code|if| and
\code{( i | a, b | c )} for \code|case| follow the same pattern, with
\code{|:} for \code|elif| and \code|ouse|. Format texts (Chapter~\ref{ch-transput}),
bits and other literals, and comments are lexical matters not shown here.
The GNU extensions \code|andth| and \code|orel| are not operators: they
bind more weakly than every operator, so that \code|a < b andth c < d|
needs no parentheses, but a chain such as \code|a andth b andth c| must
be written with parentheses, \code|a andth (b andth c)|. Modules are
described in Appendix~\ref{app-modules}.
